\documentclass[10pt,a4paper]{article}
\usepackage[margin=1.6cm]{geometry}
\usepackage{amsmath}
\usepackage{booktabs}
\usepackage{parskip}
\usepackage{titlesec}
\titlespacing*{\section}{0pt}{5pt}{1pt}
\titleformat{\section}{\normalfont\bfseries}{}{0pt}{}
\setlength{\parskip}{3pt}
\pagestyle{empty}

\begin{document}

\begin{center}
{\large\bfseries Force sensor $\rightarrow$ camera extrinsic, and the grasped-object seed}\\[1pt]
{\small Summary of values, 28 August 2026}
\end{center}

\textbf{Status.} Extrinsic resolved and validated. The grasped-object stage no longer depends on
it, or on the tool-frame question. Nothing further is needed from your side; one more recording
would help, noted at the end.

\section{The discrepancy, and its cause}
Both matrices are built from your two angles, $-135^\circ$ about $z$ and $+45^\circ$ about $x$.
They differ only in the order the two rotations are composed, and since $R_z$ and $R_x$ do not
commute that alone is the full $82.82^\circ$ between them.
\[
\underbrace{R_z(-135)\,R_x(45) =
\begin{pmatrix} -0.7071 & 0.5 & -0.5\\ -0.7071 & -0.5 & 0.5\\ 0 & 0.7071 & 0.7071\end{pmatrix}}_{\text{your sheet; \texttt{"ZXY"}, intrinsic}}
\qquad
\underbrace{R_x(45)\,R_z(-135) =
\begin{pmatrix} -0.7071 & 0.7071 & 0\\ -0.5 & -0.5 & -0.7071\\ -0.5 & -0.5 & 0.7071\end{pmatrix}}_{\text{matches recordings; \texttt{"zxy"}, extrinsic}}
\]
The sheet writes \texttt{from\_euler("Zxy", [-135,45,0])}. That string raises \texttt{ValueError}
in scipy because it mixes cases, so what ran was one or the other: \texttt{"ZXY"} reproduces your
handwritten matrix exactly, \texttt{"zxy"} gives the one matching the data. It is not a transpose
and not a flipped axis; both were tested, and no relabelling of the measurement axes alone, nor
of the camera axes alone, closes the gap.

\section{Values in use}
Applied as $p_{\text{measurement}} = T\, p_{\text{camera}}$; the inverse projects into the image.
The translation is your sheet's value unchanged, negative $z$ included.
\[
R = R_x(45)\,R_z(-135), \qquad
t = \begin{pmatrix} 0 & 120.99 & -125.65\end{pmatrix}^{\!\top}\text{mm}, \qquad
\text{grasped seed pixel } (565,\,110)
\]

\section{Results}
\begin{center}\small
\begin{tabular}{llc}
\toprule
\multicolumn{3}{l}{\textit{Extrinsic}}\\
Wrench ray lands on the contacted object & extrinsic ordering & \textbf{6 / 7 events}\\
 & intrinsic ordering & 1 / 7 events\\
 & null over 91{,}800 candidate transforms & 0.143\\
Constructions scored & 2 matrices $\times$ 3 translations $\times$ 3 placements $\times$ 2 directions & 72\\
 & best reachable from the intrinsic form & 4 / 7\\
\midrule
\multicolumn{3}{l}{\textit{Contact-object seeding, which the extrinsic unblocks}}\\
Projected contact point seeds segmentation & vs.\ 0 / 6 for the previous vision-only heuristic & \textbf{5 / 7 events}\\
Propagated track vs.\ prior hand-seeded track & mask-area ratio & 0.945\\
\midrule
\multicolumn{3}{l}{\textit{Grasped-object seeding}}\\
Seed pixel inside the held object & carrot trial & \textbf{169 / 177} (95.5\%)\\
 & cube trial & \textbf{271 / 281} (96.4\%)\\
Pixels scoring $\ge 90\%$ on both & connected region, largest component 362\,px & 398\,px\\
Validation status & fitted on both recordings & in-sample\\
\bottomrule
\end{tabular}
\end{center}

\section{Why the grasped stage no longer needs the tool frame}
On an eye-in-hand rig the arm pose cancels out of the projection, so a fixed offset in the
measurement frame is exactly a fixed pixel (verified: std $0.0$\,px across samples). The seed can
therefore be fitted directly from the recordings, with no CAD chain. For completeness, a point
104\,mm along measurement-frame $+z$ projects into the image but never lands on the held object at
any distance from $-40$ to $+300$\,mm in the carrot trial, which is why the tool-frame origin alone
could not have been the answer.

\section{Checks on your drawing, all passed}
\begin{itemize}\itemsep0pt \parskip0pt \topsep2pt
\item Handwritten matrix equals the intrinsic composition of your angles exactly.
\item Orthonormal to $3.4\times10^{-9}$, determinant $1.000$.
\item Your two translation vectors, $[174.40,\,3.3,\,0]$ global and $[0,\,120.99,\,-125.65]$ rotated,
      agree in magnitude to $0.0007$\,mm.
\item $91.72 + 12.28 = 104$\,mm, matching the CAD gripper-centre dimension.
\item Intrinsic \texttt{ZXY} $\equiv$ extrinsic \texttt{yxz}, as you noted.
\end{itemize}

\section{One thing that would help}
The grasped seed is fitted on the only two recordings that contain a carried object, so it cannot
also be held out. One further pick-and-place recording, ideally the same object from a clearly
different viewpoint, would turn it into a validated result. The pegboard trials cannot serve: the
task there is a latch moved in place, so no object is carried.

\end{document}
